\section{Conclusiones}
\label{sec:conclusiones}

El desarrollo experimental de este trabajo permitió sintetizar y caracterizar de manera comparativa tres familias distintas de puntos cuánticos de carbono (CQDs) mediante pirólisis asistida por microondas a una potencia efectiva de 770~W, demostrando que las propiedades ópticas, microestructurales y la eficiencia fotoluminiscente de los nanomateriales están fuertemente condicionadas por la naturaleza del precursor orgánico, la estequiometría de la reacción y la cinética del tratamiento térmico.

En la ruta sintética base de ácido cítrico y urea con relación molar 1:6 (4~min de pirólisis), la reacción produjo una matriz masiva monolítica y no fluorescente en estado sólido. Las micrografías SEM revelaron una estructura densa y aglomerada de bloques rocosos con porosidades generadas por el escape efervescente de gases volátiles ($\text{NH}_3$ y $\text{CO}_2$), lo que justifica el fenómeno de auto-extinción por agregación (ACQ) provocado por el empaquetamiento denso de las partículas. Por el contrario, el incremento de la relación molar a 1:10 (8~min) proporcionó el confinamiento estérico necesario para inhibir la aglomeración masiva y obtener suspensiones coloidales altamente fluorescentes.

La implementación de la ruta verde utilizando clara de huevo (ovoalbúmina) confirmó la factibilidad de sintetizar nanomateriales co-dopados con nitrógeno y azufre (N,S-CQDs) a partir de bioprecursores sostenibles. La cuantificación elemental por EDX verificó la incorporación de $0{,}46\%$ atómico de azufre ($S$), procedente de los aminoácidos cisteína y metionina de la albúmina. El estudio cinético-temporal permitió identificar una ventana cinética óptima (\textit{sweet spot}) a 1:15~minutos de pirólisis, caracterizada por un hombro de absorción a $285$~nm ($n \to \pi^*$) y un \textit{bandgap} óptico directo de $E_g = 3{,}33$~eV ($\lambda \approx 372$~nm). Superar este tiempo óptimo ($\ge 1:30$~min) desencadenó un proceso de sobre-carbonización con pérdidas de absorbancia funcional superiores al $90\%$, producto de la degradación térmica no radiativa de los fluoróforos de superficie.

Por su parte, la ruta sintética en medio ácido (ácido cítrico, urea y vinagre) demostró la capacidad del ácido acético para modificar la estructura de bandas y los estados de trampa superficiales de los puntos cuánticos. Los espectros UV-Vis revelaron un corrimiento batocrómico (desplazamiento al rojo) de $+57$~nm en la banda de absorción principal ($285$~nm $\to 342$~nm), acompañado de una elevada densidad de nitrógeno ($31{,}87\%$ at. por EDX), lo que corrobora la formación de grupos funcionales amina en la periferia de la estructura carbonosa.

Finalmente, el desarrollo e implementación del sistema de fotometría digital RGB en cámara oscura constituyó una alternativa metrológica portátil, reproducible y de bajo costo para la cuantificación óptica ante la ausencia de un espectrofluorómetro comercial. La sustracción de fondo permitió definir la luminancia neta ($Y_{\text{neta}}$), confirmando intra-fila el máximo cinético del \textit{sweet spot} ($+7{,}54$) frente a muestras sobre-carbonizadas ($+0{,}77$). Asimismo, la relación ratiométrica verde-azul ($G/B$) registró un incremento del $+170\%$ ($0{,}168 \to 0{,}453$) al pasar de la ruta verde a la ruta ácida, estableciéndose como un parámetro analítico adimensional e independiente de variaciones espaciales en la irradiancia de excitación UV.